Un verre contient \qty{\v{m}}{g} d'eau (H$_2$O). Masses molaires : $M(\text{H}) = \qty{1.008}{g/mol}$, $M(\text{O}) = \qty{16.00}{g/mol}$. Nombre d'Avogadro : \qty{6.022e23}{mol^{-1}}.

\begin{enumerate}[label=\alph*)]
    \item Calculez la masse molaire de l'eau, puis le nombre de moles d'eau dans le verre.
    \item Combien de molécules d'eau le verre contient-il ? Combien d'atomes d'hydrogène ?
    \item Quelle masse d'hydrogène l'eau du verre contient-elle ?
\end{enumerate}
